\documentclass{article}
\usepackage{amsmath,amssymb}
\title{Conditional feasibility bounds for metered learning}
\author{Emmy}
\date{}

\begin{document}
\maketitle

\section*{Source scope}
The assigned \texttt{inputs/Q.tex} is the stochastic resource allocation
mechanism of Garjani, Shokri, and Kebriaei. Its Problem Formulation defines
\(U_n=V_n(x_n)-t_n\), budget balance at equilibrium as
\(\sum_n t_n=0\), and individual rationality as \(U_n\geq V_n(0)=0\).
The assigned \texttt{inputs/P.tex} is a proposal about metering LLM decisions
in a double auction. Its Gross and net surplus section defines a real
inference expenditure \(C\) and net welfare \(G-\lambda C\). Its references
to papers called Q and P concern different source papers. The bounds below
apply only under the stated bridge assumptions.

\section*{Accounting bound}
Suppose agent \(n\)'s eventual gross gain is \(g_n\), its real inference
expenditure is \(c_n\), and \(t_n\) is a transfer among agents. Let
\(G=\sum_n g_n\), \(C=\sum_n c_n\), and let \(\lambda>0\) convert cost units
to gain units. Then the participation payoff is
\[
u_n=g_n-\lambda c_n-t_n.
\]
If total transfers balance, \(\sum_n t_n=0\), summation gives
\[
\sum_n u_n=G-\lambda C.
\]
With zero outside payoff, individual rationality for every agent requires
\(G\geq\lambda C\). Conversely, if \(G\geq\lambda C\) and arbitrary
lump-sum transfers are permitted after the outcome is known, choose any
nonnegative \(q_n\) summing to \(G-\lambda C\) and set
\(t_n=g_n-\lambda c_n-q_n\). Then transfers balance and \(u_n=q_n\geq0\).
This converse is only an accounting statement. It does not establish
incentive compatibility, Nash implementation, or information feasibility.

The need for a costless entry option already appears in Q's full valuation
class. For example, take \(\mathcal X_n=[0,1]\) and
\(V_n(x_n)=-\alpha x_n^2/2\) for every agent. These functions satisfy Q's
strong concavity and zero-at-zero assumptions, but the optimal gross value
is zero. If even the first sample call is compulsory and costs a positive
amount, total net value is negative. Balanced transfers cannot make every
agent willing to enter. This is an impossibility for a mandatory paid-sample
extension over the full type class, not for Q's free-sample mechanism.

The timing matters. In Q's Individual rationality (P4) proof, the
comparison message is \(s_n'=\operatorname{col}(0,\tilde p)\), submitted
inside the stage game after the learning process is specified. Its
nonnegative stage utility does not reimburse a cost already paid to learn
or submit that message. P's Mechanism and test section proposes a persistent
withdrawal option but charges the first scheduled decision, including a
decision to withdraw. Such withdrawal cannot serve as the costless
pre-sample entry screen required for the full valuation class in the
example above. An extension could instead offer an opt-out before the
first charged call, fund that call externally, or restrict valuations to
types with enough expected gain. A costless opt-out alone gives no useful
welfare guarantee, since all agents could exit.

For the auction schedule described in \texttt{inputs/P.tex}, Gross and net
surplus gives \(G\leq7.5\) per round. Consequently
\(\lambda C>7.5\) makes balanced transfers and universal participation
impossible for that round. This ceiling comes from P's cited auction setting;
it is not a bound on Q's resource allocation problem.

\section*{Q's growing sample schedule under a token implementation}
Q's Assumptions 1 and 2 make each \(\mathcal X_n\) compact and each
\(V_n\) continuous. Therefore
\[
M=\max_{x\in\prod_n\mathcal X_n}\sum_n V_n(x_n)<\infty.
\]
Its Theorem 4 specifies
\[
Q_i=\left\lceil\frac{C_b}{\eta^{2(i+1)}}\right\rceil,
\qquad C_b>0,\quad 0<\eta<1.
\]
Assume each of the \(NQ_i\) sample evaluations at iteration \(i\)
requires a generated-token call of real cost at least \(d>0\). This is an
additional implementation assumption: Q's original random samples need not
use language models. Put \(r=\eta^{-2}>1\). After \(I\) iterations,
\[
C_I\geq dN\sum_{i=0}^{I-1}Q_i
\geq dNC_b\sum_{i=0}^{I-1}r^{i+1}
=dNC_b\frac{r(r^I-1)}{r-1}.
\]
Suppose the allocation yields value only once after learning, the agents
have zero outside payoffs, and no external subsidy pays their cumulative
learning costs. If all transfers over the protocol balance, aggregate
participation payoff is at most \(M-C_I\). It is negative for all
sufficiently large \(I\). Thus Q's exact asymptotic schedule cannot also
provide universal ex ante individual rationality under this paid-oracle,
single-allocation interpretation. A necessary condition for an \(I\)-step
run to remain feasible is
\[
I\leq\log_r\!\left(1+\frac{M(r-1)}{dNC_b r}\right).
\]
The argument does not apply when samples are free, the cost is paid by an
outside sponsor, or repeated benefits can repay learning costs. Q's IR
theorem concerns the equilibrium payoff after learning. Q's Numerical
Simulations state \(Q_i=\lceil0.96^{i+1}\rceil=1\), which is a different
schedule from Theorem 4 and cannot validate its paid-oracle limit.

If the sample schedule and each agent's bill are fixed independently of
their messages, subtracting that bill does not change Q's stage best
responses. It changes the earlier entry decision. A behavioural mechanism
question therefore requires an endogenous choice to sample, stop, or exit,
with a finite-cost welfare or accuracy target and an explicit cost-sharing
rule. The accounting bounds alone do not provide such a mechanism.

\section*{Why the auction is not Q's allocation model}
Q assumes a separable planner objective \(\sum_nV_n(x_n)\), with strongly
concave individual valuations on convex allocation sets. P's surplus
depends on executed buyer--seller matches. For one buyer and one seller
whose calls are both needed for a gain \(h>0\), the simple binary model
\(G(x_b,x_s)=h x_bx_s\) has cross-difference
\[
G(1,1)-G(1,0)-G(0,1)+G(0,0)=h.
\]
Every additive expression \(V_b(x_b)+V_s(x_s)\) has cross-difference zero.
This example shows that applying Q's quadratic payment theorem to the
auction requires a new model of match spillovers and a relaxation of binary
participation. It does not say that all auction settings have this exact
binary payoff.

\end{document}
